% context: 9th-grade physics classwork, Unit 2 kinematics, Giancoli 2-4 and 2-5
% topic: more practice with the four constant-acceleration equations, a step easier
%        than 2-3HW: the acceleration is given in every problem but one, numbers are
%        rounder, and the elevator problem walks through average velocity first
% structure: two pages with room to work; page 1 the equations, the Vernier cart on a
%        gentler slope, and the Taipei 101 elevator; page 2 a driver's reaction and
%        braking time (a = 0, then a < 0), and a Frecciarossa braking into Milan.
%        No free fall yet.
% data: Taipei 101 elevator top speed (contexts.md #3); thinking and braking distance
%        at 30 mi/h from the UK Highway Code table, braking a from 2-3HW;
%        Frecciarossa 1000 service speed (contexts.md #2). The accelerations of the
%        elevator and train are assumed, as in 2-3HW
\documentclass[12pt, twoside]{article}
%\documentclass[12pt, twoside]{article}
\usepackage[letterpaper, margin=1in, headsep=0.2in]{geometry}
\setlength{\headheight}{0.6in}
%\usepackage[english]{babel}
\usepackage[utf8]{inputenc}
\usepackage{microtype}
\usepackage{amsmath}
\usepackage{amssymb}
%\usepackage{amsfonts}
\usepackage[nomessages]{fp} %\FPeval{\var-name}{2*sin(pi/6)}
\usepackage{siunitx} %units in math. eg 20\milli\meter
\usepackage{yhmath} % for arcs, overparenth command
\usepackage{tikz} %graphics
\usetikzlibrary{quotes, angles, arrows, arrows.meta}
\usepackage{graphicx} %consider setting \graphicspath{{images/}}
\usepackage{parskip} %no paragraph indent
\usepackage{enumitem}
\usepackage{multicol}
\usepackage{venndiagram}

\usepackage{fancyhdr}
\pagestyle{fancy}
\fancyhf{}
\renewcommand{\headrulewidth}{0pt} % disable the underline of the header
\raggedbottom
\hfuzz=2mm %suppresses overfull box warnings

\usepackage{hyperref}
\usepackage{float}

\title{Physics}
\author{Chris Huson}
\date{October 2026}

\fancyhead[LE]{\thepage}
\fancyhead[RO]{\thepage \\ Nickname: \hspace{2.25cm} \,\\ \,}
\fancyhead[LO]{La Scuola d'Italia / Huson / Physics \\* 29 October 2026}

\begin{document}

\subsubsection*{2.4 Classwork: Constant Acceleration Practice}

\textbf{Reference (provided).} \quad These hold when the acceleration $a$ is constant.

\vspace{0.1cm}
{\hspace{0.5cm}%
\renewcommand{\arraystretch}{1.8}%
\begin{tabular}{l @{\hspace{2cm}} l}
  $v = v_0 + at$ & $v^2 = v_0^2 + 2a(x - x_0)$ \\
  $x = x_0 + v_0 t + \frac{1}{2}at^2$ & $\bar{v} = \dfrac{v + v_0}{2}$
\end{tabular}}

\vspace{0.2cm}
For each problem, list what you know ($v_0$, $v$, $a$, $t$, $x - x_0$) and what you
want, then choose the equation that connects them. Show the units.

\vspace{0.2cm}

\begin{enumerate}[itemsep=0.5cm]

\item The Vernier track from 2.2 is set at a gentler slope. The cart is released from
  rest and speeds up at a constant $a = 0.200$ m/s$^2$.
  \begin{enumerate}[itemsep=0.4cm]
    \item The velocity grows by 0.200 m/s every second. Fill in the table.

      \vspace{0.1cm}
      \renewcommand{\arraystretch}{1.5}
      \begin{tabular}{| c | p{1.3cm} | p{1.3cm} | p{1.3cm} | p{1.3cm} | p{1.3cm} |}
        \hline
        $t$ (s) & \centering 0 & \centering 1.00 & \centering 2.00 &
          \centering 3.00 & \centering\arraybackslash 4.00 \\ \hline
        $v$ (m/s) & & & & & \\ \hline
      \end{tabular}
    \item How far has the cart rolled at $t = 3.00$ s?
      \vspace{3.0cm}
  \end{enumerate}

\item The elevators in Taipei 101 were the fastest in the world from 2004 to 2016,
  with a top speed of 16.8 m/s. Suppose an elevator starts from rest and speeds up
  at a constant 1.20 m/s$^2$.
  \begin{enumerate}[itemsep=0.4cm]
    \item How long does it take to reach top speed?
      \vspace{2.3cm}
    \item What is its average velocity while it speeds up?
      \vspace{2.0cm}
    \item How far does it rise while speeding up?
  \end{enumerate}

\newpage

\item A driver at 30 mi/h (13.4 m/s) sees a hazard ahead. The UK Highway Code allows
  9 m of \emph{thinking distance}, covered before the foot reaches the brake, and
  then 14 m of braking distance.
  \begin{enumerate}[itemsep=0.4cm]
    \item While the driver reacts, the car keeps going at 13.4 m/s, so $a = 0$. How
      long does the driver take to react?
      \vspace{2.8cm}
    \item In 2.3 you found that the car brakes at $a = -6.4$ m/s$^2$. How long does
      it take to stop?
      \vspace{2.8cm}
    \item What is the total time, and the total distance, from seeing the hazard to
      stopping?
      \vspace{2.0cm}
  \end{enumerate}

\item The Frecciarossa 1000 runs between Rome and Milan at 300 km/h (83.3 m/s).
  Approaching Milano Centrale, suppose it slows down at a constant 0.500 m/s$^2$,
  so \mbox{$a = -0.500$ m/s$^2$}.
  \begin{enumerate}[itemsep=0.4cm]
    \item How long does the train take to stop?
      \vspace{2.8cm}
    \item How far from the station must the driver start braking? Give your answer
      in km.
  \end{enumerate}

\end{enumerate}
\end{document}
